%% Paper 022 -- ICML-style planning + incident paper. Next Experiments leads.
%% Compile: tectonic paper_022_four_bugs_one_host_20260730.tex
\documentclass{article}
\usepackage{amsmath}\usepackage{amssymb}\usepackage{booktabs}\usepackage{graphicx}\usepackage{hyperref}
\usepackage[accepted]{icml2024}
\newcommand{\vb}{\mathrm{val\_bpb}}
\icmltitlerunning{Four Bugs, One Host}
\begin{document}
\twocolumn[
\icmltitle{Four Bugs, One Host, and a Comparison Error: Reproducing the \\ Campaign on a New Machine, Mechanistically}
\begin{icmlauthorlist}\icmlauthor{Claude Opus 5 (planning author), OPHIS}{ophis}\end{icmlauthorlist}
\icmlaffiliation{ophis}{Autoresearch reimplementation campaign}
\icmlcorrespondingauthor{Claude Opus 5 (OPHIS)}{baiyuzhu@mit.edu}
\icmlkeywords{autonomous research, reproducibility, systems debugging, sparse optimizers, Machine Learning}
\vskip 0.3in]
\printAffiliationsAndNotice{}

\begin{abstract}
The campaign's execution host became unreachable mid-session (its SSH host key rotated) and work resumed on a fresh machine with no torch environment, no cached data, and no tokenizer. Getting back to a trustworthy measurement required diagnosing and fixing four independent, mechanistically distinct failures, each traced to its root cause in source rather than patched by trial and error: (1) a data-loading path silently redirected to a leftover, unwritable mount from a different, unreachable account; (2) two breaking API changes in the `kernels` package (a publisher-trust check and a mandatory version pin) that also failed outright against a mirror endpoint lacking the relevant API surface; (3) systemd's per-user process reaping killing every backgrounded training run the instant our last SSH session closed, independent of \texttt{nohup}/\texttt{setsid}; and (4) a genuine, load-bearing bug in the adopted FlashAttention-3 backward custom ops, which silently mis-unpacked a kernel call's single-tensor return as a three-tuple -- undetected on the prior host's older kernel build, and only surfaced as a crash under this host's newer, stricter \texttt{torch.library.custom\_op} aliasing checks. We verify the fourth fix directly against the installed kernel's own source and its own default-argument convention, rather than assuming the fix from the error message alone. We then report and retract our own comparison error: having fixed the environment, we ran the \emph{current default} configuration and found it $\sim$0.004 worse than the campaign's previously reported best -- and initially treated this as an open puzzle about kernel-build numerics, when the real explanation was simpler: the previously reported best was measured under a \emph{different, non-default} configuration (\texttt{WINDOW\_PATTERN=TTTL}, \texttt{DEVICE\_BATCH\_SIZE=72}, \texttt{MATRIX\_LR=0.04}) that we were not running. A same-host, same-fix, correct-config reproduction is in flight. We close with a correction to our own prior planning paper's top-ranked proposal: it targeted the sparse n-gram optimizer path, but the champion configuration runs the \emph{dense}, \texttt{torch.compile}-fused path by default, which carries the identical full-table-decay mechanism through a different, harder-to-patch function -- and we re-score the corrected proposal accordingly.
\end{abstract}

\section{Four Bugs, Mechanistically}
Each bug is reported with its causal chain, not just its symptom, because a symptom-level fix (silence the error, retry, pin a version at random) would have left the underlying cause free to resurface under a slightly different shape -- exactly what happened between bugs 2 and 4 below.

\subsection{B1: a leftover mount from an unreachable account}
\textbf{Symptom.} \texttt{PermissionError} deep into a training run, at tokenizer load.
\textbf{Mechanism.} \texttt{lib.py} resolved its data directory with \texttt{os.path.isdir("/data")}: true on the old host (a real, writable mount), and \emph{also} true on the new host -- but there \texttt{/data} is a stale directory owned by a different, now-unreachable user, containing only a broken symlink. Existence of the directory was necessary but not sufficient; the code never checked that it was \emph{usable}.
\textbf{Fix.} Require \texttt{os.path.isfile} of the actual tokenizer file rather than directory existence. \texttt{os.path.isfile} correctly returns \texttt{False} through a broken symlink, so no filesystem write access was needed to detect the problem.

\subsection{B2: two breaking changes in a version-unpinned dependency}
\textbf{Symptom.} \texttt{ValueError: could not verify publisher trust status}, then, after a first fix, \texttt{ValueError: A kernel version or revision must be specified}.
\textbf{Mechanism.} \texttt{train.py} calls \texttt{kernels.get\_kernel(repo)} with no trust or version argument. Newer \texttt{kernels} releases require both: an explicit \texttt{trust\_remote\_code} (checked against a Hub API endpoint our mirror, \texttt{hf-mirror.com}, does not proxy -- verified: 404 on \texttt{/api/organizations/\{org\}/overview}, since \texttt{huggingface.co} itself is unreachable directly from this host) and an explicit version/revision (resolved through a \emph{different} Hub API namespace, \texttt{/api/kernels/...}, which the mirror also does not proxy -- verified: 401, not 404, suggesting an entirely unmapped path rather than a missing resource).
\textbf{Fix.} \texttt{trust\_remote\_code=True}, scoped by the fact that the call site's \texttt{repo} argument is one of exactly two hardcoded literals, never user or environment input -- and wrapped in a \texttt{try/except TypeError} fallback to a bare call, since the \texttt{kernels} version that ultimately got resolved on this host (via a later, unrelated torch pin) turned out to predate the \texttt{trust\_remote\_code} keyword entirely and raises on receiving it. The fix is therefore resilient across the version range pip may resolve, not tuned to one observed version.

\subsection{B3: systemd killed every backgrounded process on session close}
\textbf{Symptom.} Training processes, confirmed alive and progressing while an SSH session stayed open, vanished within seconds of that session closing -- despite \texttt{nohup} and \texttt{setsid}.
\textbf{Mechanism.} \texttt{loginctl show-user} reported \texttt{Linger=no} for this account. Under that setting, systemd-logind tears down a user's entire session scope -- including detached, \texttt{nohup}'d children -- the moment their last login session ends; \texttt{nohup}/\texttt{setsid} protect a process from \texttt{SIGHUP} sent by a \emph{shell}, not from a cgroup-scope teardown initiated by the \emph{service manager}. This is a common, easy-to-misdiagnose trap because the failure mode (processes die \emph{after} you finish watching them, never while you watch) looks exactly like an unrelated crash.
\textbf{Fix.} \texttt{loginctl enable-linger \$(whoami)} -- no elevated privilege required on this host. Verified by launching eight concurrent training runs, closing every SSH session, and confirming all eight were still alive and progressing on reconnect.

\subsection{B4: a custom op silently mis-unpacked a single-tensor return}
\textbf{Symptom.} \texttt{RuntimeError: bwd\_varlen ... output ... must not also be an input ... may not alias}, then, after clearing what looked like a stale compile cache, a second and different symptom: \texttt{AssertionError: wrong number of dimensions1}. Both symptoms occurred on some seeds and not others within the same launch batch, and initially looked like GPU contention or a cache race.
\textbf{Mechanism, found by reading the installed kernel's actual source rather than trial-and-error.} Both FA3 backward custom ops in \texttt{train.py} executed
\[
\texttt{dq, dk, dv, *\_ = \_fa3\_bwd\_raw(...)}
\]
assuming a \texttt{(dq, dk, dv, softmax\_d, ...)} tuple return. The installed kernel build (\texttt{varunneal/flash-attention-3}, \texttt{torch212-cxx11-cu130-x86\_64-linux}) declares \texttt{\_flash\_attn\_backward(...) -> torch.Tensor} and its body is literally \texttt{return softmax\_d} -- a single tensor. \texttt{dq}/\texttt{dk}/\texttt{dv} are \emph{out-parameters}: pre-allocated by the caller, passed as \texttt{dq=}, \texttt{dk=}, \texttt{dv=} keywords, and filled in place by the C++ call. The old code's tuple-unpacking silently sliced the single \texttt{softmax\_d} tensor along dimension 0 into three wrong-shaped tensors that additionally aliased \emph{each other} and \texttt{softmax\_d}'s storage. This is data-dependent on the document-boundary shapes for a given batch (which vary by seed, not by config), which is exactly why some seeds crashed immediately and others trained for dozens of steps first. It went undetected on the prior host because that host resolved an \emph{older} kernel build whose \texttt{\_flash\_attn\_backward} API genuinely did return a tuple -- this is an upstream API break between prebuilt kernel variants, not a bug introduced this session, and not something an isolated retry or cache clear could have fixed.
\textbf{Fix, verified against the library's own convention.} Pre-allocate \texttt{dq}/\texttt{dk}/\texttt{dv} with \texttt{torch.empty\_like} and pass them as keywords; discard the \texttt{softmax\_d} return. Before trusting this, we checked whether the kernel might require \emph{pre-zeroed} (accumulating) buffers rather than empty ones -- a subtler bug that would show up as small, non-crashing corruption rather than a clean crash. The kernel's own no-argument default path does \texttt{dq = torch.empty\_like(q) if dq is None else dq} (line 398 of the installed source): the library's own most-used code path already passes uninitialized memory, which would itself be broken if the kernel needed pre-zeroing. Our fix matches that convention exactly.
\textbf{Result.} All eight of the champion configuration's baseline seeds subsequently completed cleanly with sane, monotonically decreasing training loss -- where every seed had previously crashed within the first tens of steps.

\section{A Comparison Error, Retracted}
Having fixed the environment, we ran the scope's registered \texttt{run\_env} (\texttt{ATTN\_BACKEND=fa3}, \texttt{COMPILE\_MODE=max-autotune-no-cudagraphs}) with every other setting at \texttt{train.py}'s current defaults: \texttt{WINDOW\_PATTERN=SSSL}, \texttt{DEVICE\_BATCH\_SIZE=128}, \texttt{TOTAL\_BATCH\_SIZE=262144}, \texttt{MATRIX\_LR=0.03}. Eight seeds completed cleanly at mean $0.938618$ (sd $0.001386$), $\sim$0.0037 worse than the previously registered (and separately flagged contaminated) $0.934930$, and further worse than the campaign's previously reported best of $0.927183$.

We initially treated the gap to $0.934930$ as an open question -- candidate explanations included different low-level numerics between kernel builds, or a freshly retrained (non-identical) tokenizer, and we ruled out our own B4 fix as the cause (Section 1.4) before reporting the gap as unresolved. That caution was reasonable but the investigation stopped one step short: \textbf{the campaign's $0.927183$/$0.929541$ figures were never measured under the current default configuration at all.} They were measured under \texttt{WINDOW\_PATTERN=TTTL}, \texttt{DEVICE\_BATCH\_SIZE=72}, \texttt{TOTAL\_BATCH\_SIZE=147456}, \texttt{MATRIX\_LR=0.04} -- a configuration this campaign's own papers 017/018 had already shown outperforms the current defaults by roughly $0.007$--$0.010$, independent of anything on this host. We were comparing a corrected environment running a known-worse configuration against a number from a known-better one, and reported the resulting gap as a puzzle rather than checking which configuration actually produced the reference number. A same-host, same-fix, matched-configuration reproduction is running as this paper is written; its result belongs in the next report, not asserted here.

\section{Next Experiments}
\label{sec:next}
Per protocol, proposals are MECHANISM-only (new causal structure), and each carries a mechanism, a predicted effect with arithmetic, an exact implementation, a pre-registered numeric falsifier, a cost estimate, and 1--5 scores on novelty / provenance / validity / expected impact / reliability / feasibility-cost / falsifiability.

\subsection{E1 -- Corrected M1: lazy state decay, targeted at the actual default path}
\textbf{A correction to paper 021.} That paper's top-ranked proposal (``M1'') targeted \texttt{\_ngram\_sparse\_rmsprop\_step}, the eager, per-row optimizer path gated behind \texttt{NGRAM\_SPARSE\_GRAD=1}. The champion configuration's own \texttt{RESOLVED\_CONFIG} shows \texttt{NGRAM\_SPARSE\_GRAD=False}: the actual default path is \texttt{rmsprop\_step\_fused} (\texttt{train.py:2251--2258}), a \texttt{@torch.compile(dynamic=False, fullgraph=True)} kernel operating on the \emph{dense} gradient produced by \texttt{embedding\_dense\_backward}.

\textbf{Mechanism, re-verified against the actual default.} \texttt{rmsprop\_step\_fused} executes \texttt{exp\_avg\_sq.lerp\_(grad.square(), 1-beta2\_t)} every step -- a full-table decay over the same $\sim$2.1M$\times$384 fp32 tensor per half-table, identical in kind to what paper 021 diagnosed, just reached through the compiled dense path rather than the eager sparse one. The mechanistic argument is unchanged: at most $B{\cdot}T$ of millions of rows were touched, and the decay sweeps all of them regardless.

\textbf{Why this is harder than paper 021 assumed.} The sparse path is eager specifically because \texttt{torch.unique} has a data-dependent output shape, incompatible with \texttt{fullgraph=True} -- which is exactly why the campaign's own code comments call it out. \texttt{rmsprop\_step\_fused} being \texttt{fullgraph=True}-compiled means the lazy-decay technique's periodic renormalization (``every $K$ steps, multiply the full table by the accumulated scale and reset it'') cannot use a data-dependent Python conditional \emph{inside} the compiled region. It can still be implemented by keeping the step-scalar bookkeeping in the calling Python code (outside \texttt{torch.compile}) and dispatching to one of two compiled kernels -- an ordinary step, or a step-plus-renormalize step -- based on \texttt{step \% K}, which \texttt{torch.compile} handles cleanly since the branch is resolved before entering the graph, not inside it. This is a real, added implementation cost paper 021 did not anticipate, and downgrades feasibility relative to its original estimate.

\textbf{Predicted effect and falsifier.} Unchanged in kind from paper 021: if the decay is $\gtrsim$8\,ms of the step, removing $(K{-}1)/K$ of it clears the gate with room to spare. \textbf{Falsifier:} a CUDA-event timer around \texttt{rmsprop\_step\_fused}'s \texttt{lerp\_} call must read $\ge$8\,ms before implementation proceeds; after the change, $\le$1\,ms; paired $n{=}3$ improvement $\ge 0.003$; numerical equivalence to $<10^{-5}$ relative at the renormalization boundary.

\textbf{Scores.} novelty \textbf{4}, provenance \textbf{5} (re-verified against the actual default this time, not assumed), validity \textbf{4}, expected\_impact \textbf{4}, reliability \textbf{4}, feasibility\_cost \textbf{2} (downgraded from paper 021's implicit assumption of an eager target -- the dual-kernel dispatch adds real complexity), falsifiability \textbf{5}.

\subsection{E0 -- Round 0: an instrumented timing decomposition (prerequisite, not optional)}
\textbf{Mechanism.} One \texttt{torch.cuda.Event}-instrumented run decomposing step time into: dense n-gram gradient computation, the \texttt{rmsprop\_step\_fused} decay call specifically, backward for the rest of the model, and the FA3 varlen attention call. This resolves, cheaply and before any seed budget is spent: whether E1's target is actually the dominant cost under the \emph{default} (dense) path -- a question paper 021 answered only for the sparse path, which this configuration does not use.
\textbf{Falsifier / gate.} If the decay call reads $<$8\,ms, E1 is dead on arrival regardless of its other merits, and the next candidate (the FA3 varlen host-sync removal, deferred in paper 021 for a different reason) becomes the priority instead.
\textbf{Scores.} novelty \textbf{2}, provenance \textbf{5}, validity \textbf{5}, expected\_impact \textbf{5} (as a gate on E1's ~55 GPU-minute cost), reliability \textbf{5}, feasibility\_cost \textbf{5} (single short run), falsifiability \textbf{5}.

\subsection{Deferred, not abandoned}
Paper 021's E2 (factorized value embeddings), E3 (visit-count-decayed row LR), and E4 (the row-confidence gate) are unaffected by the dense-vs-sparse correction -- none of them touch the optimizer's decay mechanism -- and remain queued in their original form, behind E0/E1.

\section{Conclusion}
Every failure this block had a specific, checkable mechanism, and every fix was verified against that mechanism rather than against the absence of an error message. The comparison error in Section 2 is reported in the same spirit: not as a kernel-numerics mystery, but as a specific, nameable mistake -- comparing two different configurations and asking why they differ, once the actual difference is identified, is not evidence of anything except that they are different configurations. The corrected reproduction is the load-bearing result this campaign needs next, and it is already running.
\end{document}
